\subsection{HAUSDORFF SPACES}

PROPOSITION 1. Let X be a topological space. Then the following statements are equivalent:

\begin{enumerate}
\def\labelenumi{(\Alph{enumi})}
\setcounter{enumi}{7}
\tightlist
\item
  Any two distinct points of X have disjoint neighbourhoods.
\end{enumerate}

(Hi) The intersection of the closed neighbourhoods of any point of X consists of that point alone.

(Hii) The diagonal of the product space \(\mathbf{X} \times \mathbf{X}\) is a closed set.

(Hiii) For every set I, the diagonal of the product space Y = X \(^{I}\) is closed in Y.

\((\mathbf{H}^{\mathrm{iv}})\) No filter on \(\mathbf{X}\) has more than one limit point.

(Hv) If a filter \(\mathfrak{F}\) on X converges to \(x\) , then \(x\) is the only cluster point of \(\mathfrak{F}\) .

We shall prove the implications

and

\[
\begin{array}{r l} (\mathrm{H}) & \Longrightarrow (\mathrm{H} ^ {\mathrm{i}}) \Longrightarrow (\mathrm{H} ^ {\mathrm{v}}) \Longrightarrow (\mathrm{H} ^ {\mathrm{iv}}) \Longrightarrow (\mathrm{H}) \\ (\mathrm{H}) & \Longrightarrow (\mathrm{H} ^ {\mathrm{iii}}) \Longrightarrow (\mathrm{H} ^ {\mathrm{ii}}) \Longrightarrow (\mathrm{H}). \end{array}
\]

\((\mathrm{H}) \Longrightarrow (\mathrm{H}^{\mathrm{i}}): \text{ If } x \neq y \text{ there is an open neighbourhood } \mathrm{U} \text{ of } x \text{ and an open neighbourhood } \mathrm{V} \text{ of } y \text{ such that } \mathrm{U} \cap \mathrm{V} = \varnothing; \text{ hence } y \notin \overline{\mathrm{U}}.\)

\((\mathbf{H}^{\mathrm{i}}) \Longrightarrow (\mathbf{H}^{\mathrm{v}}): \text{ Let } y \neq x; \text{ then there is a closed neighbourhood } V \text{ of } x \text{ such that } y \notin V, \text{ and by hypothesis there exists } M \in \mathfrak{F} \text{ such that } M \subset V; \text{ thus } M \cap \mathbb{C}V = \emptyset. \text{ But } \mathbb{C}V \text{ is a neighbourhood of } y; \text{ hence } y \text{ is not a cluster point of } \mathfrak{F}.\)

\((\mathbf{H}^{\mathrm{v}}) \Rightarrow (\mathbf{H}^{\mathrm{iv}})\) : Clear, since every limit point of a filter is also a cluster point.

\((\mathrm{H}^{\mathrm{iv}}) \Longrightarrow (\mathrm{H})\) : Suppose \(x \neq y\) and that every neighbourhood \(V\) of \(x\) meets every neighbourhood \(W\) of \(y\) . Then the sets \(V \cap W\) form a base of a filter which has both \(x\) and \(y\) as limit points, which is contrary to hypothesis.

\((\mathrm{H}) \Longrightarrow (\mathrm{H}^{\mathrm{iii}})\) : Let \((x) = (x_{\iota})\) be a point of \(\mathbf{X}^{\mathbf{I}}\) which does not belong to the diagonal \(\Delta\) . Then there are at least two indices \(\lambda, \mu\) such that \(x_{\lambda} \neq x_{\mu}\) . Let \(\mathbf{V}_{\lambda}\) (resp. \(\mathbf{V}_{\mu}\) ) be a neighbourhood of \(x_{\lambda}\) (resp. \(x_{\mu}\) ) in \(\mathbf{X}\) , such that \(\mathbf{V}_{\lambda} \cap \mathbf{V}_{\mu} = \emptyset\) ; then the set \(W = V_{\lambda} \times V_{\mu} \times \prod_{\iota \neq \lambda, \mu} X_{\iota}\) (where \(X_{\iota} = X\) if \(\iota \neq \lambda, \mu\) ) is a neighbourhood of \(x\) in \(X^{\mathbf{I}}\) ( \(\S 4, no.1\) ) which does not meet \(\Delta\) . Hence \(\Delta\) is closed in \(X^{\mathbf{I}}\) .

\((\mathbf{H}^{\mathrm{iii}}) \Rightarrow (\mathbf{H}^{\mathrm{ii}}): \text{Obvious.}\)

\((H^{ii}) \Longrightarrow (H)\) : If \(x \neq y\) then \((x, y) \in X \times X\) is not in the diagonal \(\Delta\) , hence (§ 4, no. 1) there is a neighbourhood \(V\) of \(x\) and a neighbourhood \(W\) of \(y\) in \(X\) such that \((V \times W) \cap \Delta = \emptyset\) , which means that

\[
\mathbf {V} \cap \mathbf {W} = \varnothing .
\]

DEFINITION 1. A topological space satisfying the conditions of Proposition 1 is called a Hausdorff or separated space; the topology of such a space is said to be a Hausdorff topology.

Axiom (H) is Hausdorff's axiom.

Examples. Any discrete space is Hausdorff. The rational line Q is Hausdorff, for if x, y are two rational numbers such that x \textless{} y then there is a rational number z such that x \textless{} z \textless{} y, and the neighbourhoods {]}\(\leftarrow\) , z{[} of x and {]}z, \(\rightarrow\) {[} of y do not intersect.

A set X having at least two points and carrying the coarsest topology (§ 2, no. 2) is not a Hausdorff space.

Let \(f: X \to Y\) be a mapping of a set \(X\) into a Hausdorff space \(Y\) ; then it follows immediately from Proposition 1 that \(f\) can have at most one limit with respect to a filter \(\mathfrak{F}\) on \(X\) , and that if \(f\) has \(y\) as a limit with respect to \(\mathfrak{F}\) , then \(y\) is the only cluster point of \(f\) with respect to \(\mathfrak{F}\) .

PROPOSITION 2. Let \(f, g\) be two continuous mappings of a topological space \(\mathbf{X}\) into a Hausdorff space \(\mathbf{Y}\) ; then the set of all \(x \in \mathbf{X}\) such that \(f(x) = g(x)\) is closed in \(\mathbf{X}\) .

For this set is the inverse image of the diagonal of \(\mathbf{Y} \times \mathbf{Y}\) under the mapping \(x \to (f(x), g(x))\) , which is continuous ( \(\S 4\) , no. 1, Proposition 1). The result therefore follows from (H \(^{ii}\) ) and \(\S 2\) , no. 1, Theorem 1.

COROLLARY I (Principle of extension of identities). Let \(f, g\) be two continuous mappings of a topological space \(X\) into a Hausdorff space \(Y\) . If \(f(x) = g(x)\) at all points of a dense subset of \(X\) , then \(f = g\) .

In other words, a continuous map of X into Y (Hausdorff) is uniquely determined by its values at all points of a dense subset of X.

COROLLARY 2. If \(f\) is a continuous mapping of a topological space \(\mathbf{X}\) into a Hausdorff space \(\mathbf{Y}\) , then the graph of \(f\) is closed in \(\mathbf{X} \times \mathbf{Y}\) .

For this graph is the set of all \((x, y) \in \mathbf{X} \times \mathbf{Y}\) such that \(f(x) = y\) , and the two mappings \((x, y) \to y\) and \((x, y) \to f(x)\) are continuous.

PROPOSITION 3. Let \((x_{i})_{1\leqslant i\leqslant n}\) be a finite family of distinct points of a Hausdorff space \(\mathbf{X}\) ; then each \(x_{i}\) has a neighbourhood \(\mathbf{V}_{i}\) in \(\mathbf{X}\) such that the \(\mathbf{V}_{i}\) ( \(1\leqslant i\leqslant n\) ) are mutually disjoint.

The proof is by induction on \(n\) : the case \(n = 2\) is just the axiom (H). Let then \(W_i\) ( \(1 \leqslant i \leqslant n - 1\) ) be a neighbourhood of \(x_i\) such that the \(W_i\) are mutually disjoint. On the other hand, for \(1 \leqslant i \leqslant n - 1\) there is a neighbourhood \(T_i\) of \(x_i\) and a neighbourhood \(U_i\) of \(x_n\) which do not intersect. If we take \(V_i\) to be \(W_i \cap T_i\) for \(1 \leqslant i \leqslant n - 1\) and

\[
\mathrm{V} _ {n} = \bigcap_ {i = 1} ^ {n - 1} \mathrm{U} _ {i},
\]

the conditions of the proposition are satisfied.

COROLLARY. Every finite Hausdorff space is discrete.

PROPOSITION 4. Every finite subset of a Hausdorff space is closed.

For every subset consisting of a single point is closed by reason of axiom (H \(^{i}\) ).

PROPOSITION 5. Let X be a topological space and suppose that for each pair of distinct points x, y of X there is a continuous mapping f of X into a Hausdorff space \(X'\) such that \(f(x) \neq f(y)\) . Then X is Hausdorff.

Let \(V'\) and \(W'\) be disjoint neighbourhoods of \(f(x)\) and \(f(y)\) respectively in \(X'\) ; then \(\overline{f}^1(V')\) and \(\overline{f}^1(W')\) are disjoint neighbourhoods of \(x\) and \(y\) respectively in \(X\) .

COROLLARY. Every topology which is finer than a Hausdorff topology is Hausdorff.

